\documentclass[10pt]{article}

\usepackage[margin=1in]{geometry}
\usepackage{amsmath,amssymb}
\usepackage{booktabs}
\usepackage{microtype}
\usepackage[hidelinks]{hyperref}
\usepackage{xcolor}

\newcommand{\systemname}{GEODE}
\newcommand{\ba}{\operatorname{BA}}
\newcommand{\utility}{U}
\setlength{\emergencystretch}{2em}

\title{Adaptation Utility Reveals Why Qualified Open-World Recognition
Stages Fail to Compose}
\author{
  Moazzam Abdullah Khan\\
  \texttt{moazzam2k@yahoo.com}
}
\date{}

\begin{document}
\maketitle

\begin{abstract}
Open-world recognition is usually evaluated as a sequence of separate
problems: reject unfamiliar inputs, cluster rejected examples, request labels,
add classes, and route future inputs. Strong stage-level results do not
guarantee that these stages compose into a useful system. We study this gap in
a preregistered, human-review-budgeted loop over frozen 384-dimensional DINOv2
features. The protocol uses one primary endpoint: post-integration balanced
accuracy on known and discovered classes at 50 reviewed labels per episode,
with known-class regression and remaining-unknown recall as safety constraints.

The parent study first established that several stages could pass in
isolation. A rank-32 class-conditional Gaussian achieved 92.08\% known
coverage, 86.17\% unknown recall, 0.9561 AUROC, and 98.67\% review precision.
HDBSCAN achieved 100\% group recall and precision in its stage-wise
evaluation, and confirmed rank-16 class insertion improved isolated target
success by 43.33 percentage points. Nevertheless, grouped end-to-end loops
integrated 0 of 3 confirmable classes.

We then optimize review selection for adaptation utility rather than purity.
Across nine seed-by-arrival cells, deterministic utility selection beats
density-core selection in all nine cells by 3.593 balanced-accuracy points on
average, with a paired 95\% bootstrap interval of [2.794, 4.382]. This is a
real and consistent signal, but it misses the preregistered 5-point minimum and
violates the 2-point remaining-unknown-recall band in six cells. A separate
local-residual study preserves 100\% of unaffected predictions, repairing an
earlier scope-leakage failure, but adds at most 0.204 points of utility and
also fails unknown-recall safety. The final registered outcome is therefore
negative: purity-selected discovery outputs are insufficient statistics for
safe adaptation, while improving representativeness helps but does not by
itself qualify an open-world lifecycle. The study contributes an auditable
episode protocol, typed stage interfaces, fail-closed lineage contracts, exact
rollback, and artifact-only reproduction of the negative result.
\end{abstract}

\noindent\textbf{Keywords:} open-world recognition, open-set recognition,
novel-class discovery, human-in-the-loop learning, continual learning,
adaptation utility, transactional machine learning, negative results

\section{Introduction}

A deployed classifier eventually encounters examples outside its original
class inventory. Open-world recognition requires the system to recognize known
classes, reject unknowns, and incrementally add newly supervised classes
\cite{bendale2015}. In an operational system, however, rejection is only the
first step. Rejected evidence must be buffered; recurring structure must be
separated from corruption; candidate groups must be reviewed; semantic labels
must be supplied by an authorized actor; updates must be validated; and every
publication must remain auditable and reversible.

The literature contains mature mechanisms for individual parts of this loop.
OpenMax and the Extreme Value Machine calibrate unfamiliarity
\cite{bendale2016,rudd2018}. Stream-mining systems buffer and organize emerging
classes \cite{masud2011,haque2016sand,haque2016echo,faria2016,mu2017}.
Generalized and continual category-discovery methods identify novel groups in
modern representation spaces \cite{vaze2022,cao2022,zhang2022,zhao2023}.
Expert Gate and modular model systems route among specialists
\cite{aljundi2017,gururangan2021,li2022}. Yet independently successful stages
may emit the wrong sufficient statistics for the next stage.

This paper asks:

\begin{quote}
Can a review-budgeted discovery, confirmation, adaptation, and routing loop
improve post-integration performance while preserving known-class accuracy,
remaining-unknown recall, and transactional safety?
\end{quote}

We investigate this question in \systemname, an experimental framework for
explicit geometric models and auditable model lifecycles. We do not assume that
signed-distance geometry is necessary or competitively optimal. In the same
frozen representation, an earlier weighted affine head reached 91.73\%
balanced accuracy, compared with 96.77\% for an RBF SVM. The acceptance
mechanism is therefore treated as substitutable, and a low-rank Gaussian is
the retained primary head.

Our contributions are:

\begin{enumerate}
  \item \textbf{A registered end-to-end endpoint.} We evaluate adaptation
  utility at a fixed human-review budget instead of promoting systems through
  purity, AUROC, or coverage alone.

  \item \textbf{A quantified interface failure.} Dense core sets cover only
  34.0--39.1\% of the rejected region under a nearest-neighbor support
  diagnostic, whereas boundary-inclusive sets cover 99.37--99.78\%.

  \item \textbf{A consistent but insufficient selection signal.} Utility
  selection improves over core selection in 9/9 cells, but the gain is smaller
  than registered and does not preserve remaining-unknown recall.

  \item \textbf{A transactional harness.} Typed interface contracts,
  immutable lineage, human-confirmation gates, exact rollback, deterministic
  replay, and artifact-only conclusion verification separate predictive
  failure from systems failure.
\end{enumerate}

The central result is negative but actionable. Optimizing discovery for
review purity does not provide sufficient support for adaptation. Optimizing
support diversity improves adaptation, but safe open-world learning remains a
joint constrained problem rather than a sequence of independently optimized
stages.

\section{Related Work}

\subsection{Open-set and open-world recognition}

Bendale and Boult formalized open-world recognition as known-class
classification with explicit rejection and incremental class addition
\cite{bendale2015}. OpenMax calibrates activation vectors using extreme-value
theory \cite{bendale2016}, while the Extreme Value Machine models
class-conditional margin tails and supports incremental updates
\cite{rudd2018}. Feature-space alternatives include maximum posterior
probability \cite{hendrycks2017}, Mahalanobis or class-conditional density
scores \cite{lee2018,mukhoti2023}, deep nearest neighbors \cite{sun2022}, and
distance-aware discriminative models \cite{vanamersfoort2020,liu2020}.

These methods mainly answer whether an input lies outside learned support. Our
primary question starts after rejection: whether selected rejected examples
contain enough information to create a useful and safe child model.

\subsection{Streaming novel-class detection}

Data-stream systems operationalized reject-buffer-group-update loops before
the modern deep-representation era. ECSMiner detects novel classes under
concept drift \cite{masud2011}; SAND and ECHO address limited labels and
evolving concepts \cite{haque2016sand,haque2016echo}; MINAS maintains
micro-clusters and recurring novelty \cite{faria2016}; and SENCForest targets
emerging classes \cite{mu2017}. These systems motivate our persistent buffer
and delayed-confirmation protocol. We add immutable transaction lineage and
measure downstream adaptation utility directly.

\subsection{Category discovery}

Generalized category discovery clusters unlabeled data containing both known
and unknown categories, often over strong self-supervised features
\cite{vaze2022}. ORCA jointly addresses seen and novel classes in open-world
semi-supervised learning \cite{cao2022}. Grow and Merge and incremental GCD
extend discovery across time \cite{zhang2022,zhao2023}. Our study holds the
DINOv2 representation \cite{oquab2024} fixed so that review-set selection,
rather than representation training, is the manipulated variable.

DBSCAN \cite{ester1996}, HDBSCAN \cite{campello2013}, and FINCH
\cite{sarfraz2019} offer distinct grouping mechanisms. Cluster quality is
usually measured through coherence, purity, or recovery. We show that a pure
cluster core can still be an inadequate adaptation set.

\subsection{Expert routing and explicit models}

Expert Gate uses reconstruction error to route among task experts
\cite{aljundi2017}. More recent modular systems compose or route among
independently trained experts \cite{gururangan2021,li2022,huang2024}.
Our protocol distinguishes compatibility from empirical support. A model
fingerprint determines whether a bundle may be called; a learned support
profile may only propose a shortlist. Stale or low-confidence routes fall back
to exhaustive compatible-model evaluation.

Explicit geometry offers inspectable and editable artifacts
\cite{requicha1980,park2019}, while low-rank Gaussian class models are related
to probabilistic PCA and mixtures of factor analyzers
\cite{tipping1999,ghahramani1996}. Editability is valuable for governance, but
does not imply predictive superiority. We therefore report geometric heads as
ablations rather than premises.

\section{Problem Formulation}

Let $f(x)\in\mathbb{R}^d$ be a frozen representation. Before episode $e$, the
parent model recognizes class inventory $\mathcal{K}_e$. Review and
confirmation may add a discovered class set $\mathcal{D}_e$. The registered
episode utility is

\begin{equation}
\utility_e =
\ba_{\mathcal{K}_e\cup\mathcal{D}_e}(\text{child})
-
\ba_{\mathcal{K}_e\cup\mathcal{D}_e}(\text{parent}).
\label{eq:utility}
\end{equation}

Rejected or unavailable discovered labels count as incorrect for the parent.
This prevents a no-update system from appearing competitive by excluding the
new class from evaluation.

\subsection{Fixed resources and safety}

Each episode receives exactly $B=50$ reviewed labels. Unspent labels do not
transfer. The representation and episode partitions are immutable. A candidate
must keep known-class regression below one percentage point and preserve
remaining-unknown recall within two points of its parent. Semantic publication
requires linked human confirmation. Replay, graph validation, rollback,
lineage, and fallback checks must all pass.

\subsection{Statistical advancement}

The pivotal review-selection hypothesis is preregistered. Utility-selected
review advances only if it:

\begin{enumerate}
  \item exceeds density-core selection by at least 5.0 balanced-accuracy points
  in mean $\utility_e$;
  \item has a paired 95\% bootstrap lower bound above zero;
  \item improves at least 7 of 9 seed-by-episode cells; and
  \item satisfies the review-budget, safety, and transaction contracts.
\end{enumerate}

The first criterion is an effect-size requirement, not a significance test.
Passing the interval and cell-count criteria is insufficient if the practical
or safety thresholds fail.

\section{System and Protocol}

\subsection{Data, representation, and episodes}

The realized study is a bounded CIFAR-10 proxy using immutable
384-dimensional CLS-token features from DINOv2-small. Each frozen seed has
10,000 training and 1,000 development examples. Seeds are 11, 23, and 37.
The v8 replay begins with classes 0--5 and introduces classes 6, 7, and 8 in
order, yielding nine seed-by-arrival cells. Class 9 remains unknown. No
untouched final-label schedule is opened.

Each class is deterministically partitioned into geometry fitting, anchor,
review-candidate, and proxy-validation subsets. The validation subset may
score a frozen candidate review set but never contributes adaptation support.
Every episode records class order, partition hashes, support IDs, confirmation
IDs, thresholds, parent and child hashes, utility, safety metrics, and rollback
evidence.

\subsection{Stage-interface contracts}

We register five adjacent interfaces:

\begin{enumerate}
  \item rejector $\rightarrow$ buffer;
  \item buffer $\rightarrow$ clusterer;
  \item clusterer $\rightarrow$ review selector;
  \item review $\rightarrow$ adapter; and
  \item adapter $\rightarrow$ router.
\end{enumerate}

Seven versioned schemas represent the utility endpoint, episode replay,
interface audits, threshold transfer, review selection, local residual scope,
and integrated routing. The harness aborts on partition leakage, stale class
order, threshold-lineage mismatch, review-budget overflow, missing
confirmation, or rollback-parent drift.

\subsection{Acceptance and threshold transfer}

The retained acceptance head is a rank-32 class-conditional Gaussian. For
center $\mu$, orthonormal tangent basis $V$, tangent variances $\lambda_j$,
and residual variance $\sigma^2$,

\begin{equation}
q(x)=
\sum_j\frac{(v_j^\top(x-\mu))^2}{\lambda_j}
+
\frac{\lVert(I-VV^\top)(x-\mu)\rVert_2^2}{\sigma^2}.
\end{equation}

The class log likelihood includes the corresponding log determinant, and
novelty is $-\max_k\log p(x\mid k)$. We compare a frozen parent threshold, a
class-count heuristic, a global anchor quantile, and per-class anchor
thresholds with a global fail-closed ceiling. Stale anchor lineage restores the
parent threshold.

\subsection{Review-selection arms}

All eligible arms receive the same candidate pool and exactly 50 labels:

\begin{description}
  \item[Core-selected.] Samples nearest the candidate-group center.
  \item[Boundary-inclusive.] A deterministic mix of core examples, modes, and
  parent-boundary-distance quantiles.
  \item[Coverage-selected.] Farthest-first $k$-center selection.
  \item[Utility-selected.] Deterministic candidate subsets are ranked by a
  cheap rank-3 proxy adapter on a disjoint validation partition; the production
  adapter remains rank 16.
  \item[Random stratified.] Seeded sampling across two candidate modes.
  \item[Review-everything.] All 300 candidate examples; diagnostic only and
  ineligible for equal-budget advancement.
\end{description}

All selected samples belong to the arriving class in this controlled
sufficient-statistics study. Purity is therefore fixed at 100\%; the
experiment isolates representativeness.

\subsection{Transactional adaptation and routing}

The production adapter adds one confirmed rank-16 low-rank component and
recalibrates the global anchor threshold. Each transaction appends exactly one
class, links a persistent review and confirmation, and must restore the exact
parent bundle under rollback. Routing remains exhaustive because no empirical
router passed the parent study's authoritative gate.

\subsection{Local residual branch}

Two preregistered residuals are evaluated independently after the pivotal
selection experiment: a local target-class temperature and a bounded
class-local affine correction. Before fitting, the parent alone freezes the
affected region using selected support, a bounded nearest-neighbor radius, and
component responsibility. Residual fusion is applied only inside that mask.
An attempt requires at least 5 points of utility improvement, 99.9\%
unaffected-prediction preservation on every seed, and the same known and
unknown safety bounds.

\section{Results}

\subsection{Parent stage-wise evidence}

Table~\ref{tab:parent} summarizes the stage-level mechanisms inherited from the
v7 parent study. These results motivated composition but did not establish it.

\begin{table}[t]
\centering
\caption{Parent stage-wise results. Values are three-seed means.}
\label{tab:parent}
\begin{tabular}{llr}
\toprule
Stage & Retained mechanism & Result \\
\midrule
Rejection & Rank-32 Gaussian & 86.17\% unknown recall \\
           &                  & 92.08\% known coverage \\
           &                  & 98.67\% review precision \\
Discovery & HDBSCAN          & 100\% group recall/precision \\
Discovery & FINCH            & 83.33\% recall; 86.39\% precision \\
Adaptation & Confirmed rank-16 insertion & $+43.33$ points target success \\
Routing & Best shadow profile & 99.57\% winner inclusion \\
\bottomrule
\end{tabular}
\end{table}

Despite these stage-level results, HDBSCAN and FINCH integrated 0/3
confirmable classes in the v7 grouped loop. Reviewing every rejection
integrated 3/3 but provided no review reduction. This composition failure
motivated the v8 endpoint.

\subsection{Threshold transfer}

Frozen, class-count, and global anchor-quantile thresholds all satisfy the
registered bands on seeds 11, 23, and 37. We retain the global anchor rule
because it directly recalibrates after class-order change. Mean known-class
regression and mean remaining-unknown-recall change are both 0.00 points; mean
remaining-unknown recall is 79.33\%. Per-class fail-closed thresholds increase
unknown rejection but cause 1.875 points of mean known regression, exceeding
the one-point ceiling.

\begin{table}[t]
\centering
\caption{Threshold transfer over three development seeds.}
\label{tab:threshold}
\begin{tabular}{lrrc}
\toprule
Rule & Mean known drop & Mean unknown drop & Pass \\
\midrule
Frozen parent & 0.000 & 0.000 & Yes \\
Class-count heuristic & 0.000 & 0.000 & Yes \\
Global anchor quantile & 0.000 & 0.000 & Yes \\
Per-class anchor & 1.875 & $-8.667$ & No \\
\bottomrule
\end{tabular}
\end{table}

\subsection{Core sets are geometrically unrepresentative}

Table~\ref{tab:coverage} compares 50 core examples with 50
boundary-inclusive examples. Both cover the two coarse $k$-means modes, so
mode count alone does not expose the defect. Nearest-neighbor support and
covariance do. Core sets cover only 34.0--39.1\% of the rejected candidate
region; boundary-inclusive sets cover more than 99\% and approximately recover
the full covariance trace.

\begin{table}[t]
\centering
\caption{Review-set representativeness diagnostics.}
\label{tab:coverage}
\begin{tabular}{crrrr}
\toprule
Seed & Core NN cov. & Boundary NN cov. & Core cov. ratio & Boundary cov. ratio \\
\midrule
11 & 34.04\% & 99.68\% & 0.663 & 1.078 \\
23 & 35.47\% & 99.78\% & 0.658 & 1.096 \\
37 & 39.14\% & 99.37\% & 0.679 & 1.068 \\
\bottomrule
\end{tabular}
\end{table}

In the M46 proxy adapter, boundary inclusion adds 0.259 balanced-accuracy
points over core selection on average. The effect is positive on two of three
seeds but is much smaller than the pivotal M47 threshold. The interface audit
ranks changed-region fusion scope, full normalized router scores, and boundary
examples as the three measured mismatches.

\subsection{Equal-budget adaptation utility}

Table~\ref{tab:utility} reports mean episode utility across nine cells.
Coverage selection is the strongest eligible deterministic arm at 10.62
points. Utility selection reaches 9.96 points, compared with 6.36 for core
selection. Review-everything reaches 11.26 points but uses 300 labels per
episode and is ineligible.

\begin{table}[t]
\centering
\caption{Mean adaptation utility and review cost across nine cells.}
\label{tab:utility}
\begin{tabular}{lrrr}
\toprule
Review policy & Mean $\utility_e$ & Mean labels & Mean remaining-unknown recall \\
\midrule
Core-selected & 6.36 & 50 & 59.72\% \\
Boundary-inclusive & 9.71 & 50 & 75.98\% \\
Coverage-selected & \textbf{10.62} & 50 & \textbf{82.50\%} \\
Utility-selected & 9.96 & 50 & 78.15\% \\
Random stratified & 9.70 & 50 & 75.41\% \\
Review-everything & 11.26 & 300 & 86.13\% \\
\bottomrule
\end{tabular}
\end{table}

Utility selection beats core selection in 9/9 cells. The mean paired
difference is 3.593 points, and the 10,000-resample paired percentile
bootstrap interval is [2.794, 4.382]. Thus the sign is stable and the interval
excludes zero. Nevertheless, the gain is below the registered 5-point practical
threshold.

The safety conjunction also fails. Every utility-selected transaction consumes
exactly 50 reviews, links confirmation, appends one class, replays, and rolls
back exactly. Mean known regression is negative, indicating improvement rather
than degradation. However, six of nine cells lose more than two points of
remaining-unknown recall. The registered pivotal gate therefore fails and
blocks learned selection and final E12 qualification.

\begin{table}[t]
\centering
\caption{Preregistered pivotal gate.}
\label{tab:gate}
\begin{tabular}{lrc}
\toprule
Operand & Observed & Pass \\
\midrule
Mean gain over core $\geq 5.0$ points & 3.593 & No \\
Paired 95\% lower bound $>0$ & 2.794 & Yes \\
Positive cells $\geq 7/9$ & 9/9 & Yes \\
Budget, safety, and transaction conjunction & 6 unknown-recall violations & No \\
\bottomrule
\end{tabular}
\end{table}

\subsection{Locality can be repaired without creating utility}

The parent A3 study preserved only 84.2\% of unaffected predictions because
an ostensibly local edit entered a global fusion path. Parent-only region masks
and explicit fusion gating raise unaffected preservation to 100\% for both M49
attempts. Maximum known regression is only 0.125 points.

The residuals do not qualify. Local target temperature adds 0.000 points of
mean utility. The bounded affine correction adds 0.204 points, far below the
5-point gate. Both observe a maximum remaining-unknown-recall drop of 4.333
points. No residual is retained.

\section{Discussion}

\subsection{Purity is not an adaptation endpoint}

The parent study appears strong when read stage by stage: high unknown recall,
near-perfect review precision, perfect HDBSCAN recovery, and large isolated
new-class gains. The integrated failure and v8 selection results show why this
reading is misleading. A dense core can be semantically pure while discarding
low-density modes and boundary variation needed by the child model.

The 9/9 utility gain over core selection is evidence for the sufficient-
statistics hypothesis. The failure of the 5-point and safety gates is equally
important. Representativeness is necessary in this harness, but not sufficient.
Future objectives must jointly optimize adaptation utility and preservation of
the remaining unknown space.

\subsection{Why the proxy selector is not the winner}

Coverage-selected review outperforms utility-selected review. The utility
selector ranks subsets with a cheap rank-3 proxy but deploys a rank-16
production adapter. This creates proxy regret: validation utility under the
cheap model does not perfectly rank production outcomes. Training a learned
selector might reduce regret, but the preregistered kill switch correctly
prohibits that continuation after the pivotal endpoint fails.

\subsection{Unknown recall is a lifecycle state variable}

Adding a class changes the score competition and the meaning of a calibrated
threshold. Global anchor recalibration transfers safely in the isolated M46
experiment, yet six M47 cells still exceed the remaining-unknown band. This
shows that threshold transfer cannot be evaluated independently from support
selection and cumulative class addition. Unknown recall must be measured after
each transaction, not inherited from the rejection benchmark.

\subsection{Locality and rollback answer different questions}

Rollback proves that the parent can be restored. It does not prove that the
child update is local. The M49 result demonstrates a direct repair:
parent-defined masks plus explicitly gated fusion preserve every unaffected
prediction. The absence of utility gain prevents a stronger claim, but the
mechanism clarifies how edit scope should be enforced.

\subsection{Negative results and stopping rules}

The final outcome is not ``no signal.'' Utility selection has a positive,
consistent effect. The outcome is negative because the effect is smaller than
the registered practical threshold and violates a safety co-primary.
Publishing only the positive interval would overstate the evidence. The
stopping rule prevents post-hoc learner tuning, head substitution, or untouched
schedule access from converting a failed primary endpoint into a new claim.

\section{Reproducibility}

All milestone configurations, schemas, selected sample IDs, parent and child
lineage, and result artifacts are stored in the project repository. The final
artifact-only command is:

\begin{verbatim}
.\.venv\Scripts\python.exe -m experiments.tier1.eval_v8_final_replay
\end{verbatim}

The verifier checks four immutable milestone indexes, 11 indexed artifacts,
and six conclusion operands. Two runs must be byte-identical. The verifier
loads no training data and opens no final labels. The complete repository test
gate contains 374 tests.

For an archival release, the authors should insert the public repository URL
and release DOI here:
\href{https://example.com/repository}{repository URL placeholder}.
The arXiv source package should include this \texttt{.tex} file; no generated
evidence or benchmark data are required to compile the manuscript.

\section{Successor Geometric-Support Studies}

After the lifecycle outcome was frozen, two separately registered studies
examined whether class support in the same frozen features was better described
by boundaries or higher-codimension manifolds. The v9 study found no calibrated
shell occupancy around the existing zero-level sets. Rank-16 and rank-32 affine
tubes improved seed-11 known balanced accuracy by 1.375 and 1.625 points, but
accepted every 8x tangent-extrapolation probe because residual and tangent terms
were not comparably scaled.

V10 corrected this defect with dimensionless residuals, rank-normalized tangent
overshoot, and calibration-only safety selection. Controlled ambient-64 tests
recovered ranks 8, 16, and 32 and rejected all 8x probes. On the frozen real
features, however, 16 of 18 registered cells were infeasible during calibration.
best feasible cell passed all open-space gates but improved known balanced
accuracy by only 0.875 points, below the registered 1-point gate. Thus the
successor program also ended in Outcome D: controlled identifiability and a
small residual signal did not qualify a practical support model or reopen
lifecycle evaluation.

\section{Limitations}

\paragraph{Proxy scale.}
The realized study uses CIFAR-10 classes and frozen DINOv2 features. It does
not establish performance on large semantic inventories, natural long-running
streams, or domain-scale routing.

\paragraph{Development-only evidence.}
The untouched M50 schedule remains sealed because the M47 kill switch fired.
The result is preregistered development evidence, not an independent final-test
confirmation.

\paragraph{Simulated review.}
A deterministic oracle models human confirmation. Real reviewers introduce
disagreement, delay, expertise variation, and context-dependent cost.

\paragraph{Controlled candidate purity.}
The v8 sufficient-statistics experiment uses candidates from the true arriving
class to isolate review-set geometry. Real discovery also introduces false
members and fragmented groups.

\paragraph{Limited selectors and adapters.}
The deterministic utility selector enumerates registered subset families and
uses a rank-3 proxy. The result does not rule out better constrained selectors,
but the registered program prohibits training them after the pivotal failure.

\paragraph{No head-substitution conclusion.}
The same-harness geometric and kNN substitutions are not run after the binding
kill switch. Earlier stage-level comparisons cannot be relabeled as
adaptation-utility evidence.

\paragraph{No universal novelty claim.}
The prior-art audit cannot rule out patents, non-English work, inaccessible
publications, or closed industrial systems.

\section{Ethical and Deployment Considerations}

Open-world systems can assign semantic meaning to coherent but irrelevant,
sensitive, or harmful clusters. This study prohibits autonomous semantic
publication and requires linked confirmation, immutable provenance, and exact
rollback. These controls reduce but do not eliminate risk. Real deployment
should additionally support reviewer disagreement, appeal and deletion
procedures, privacy-aware retention, subgroup performance audits, monitoring
for feedback loops, and explicit authority boundaries.

The experiments use public benchmark representations and simulated review.
They do not involve human subjects or sensitive personal data. The architecture
should not be interpreted as authorization for high-stakes use.

\section{Conclusion}

We evaluated an open-world learning loop through adaptation utility rather than
stage-wise proxies. Strong rejection, grouping, and isolated adaptation results
did not compose. Review sets selected for density-core purity were
geometrically unrepresentative. Boundary- and coverage-aware selection improved
post-integration balanced accuracy consistently, but the registered
utility-selected arm gained only 3.593 points rather than 5 and failed the
remaining-unknown safety constraint.

The broader methodological lesson is that open-world recognition should be
evaluated as a constrained lifecycle. Rejection recall, cluster purity,
adaptation accuracy, routing inclusion, locality, and rollback can all look
strong independently while the complete system remains unqualified. The
interfaces between stages---and the information each consumer actually
needs---must be treated as first-class experimental objects.

The successor support studies reinforce this lesson. Correcting a geometric
score's units and proving synthetic identifiability were necessary but not
sufficient: system-level component masking and a missed predictive threshold
prevented the bounded support model from advancing.

\begin{thebibliography}{32}

\bibitem{bendale2015}
A. Bendale and T. Boult.
\newblock Towards open world recognition.
\newblock In \emph{CVPR}, 2015.

\bibitem{bendale2016}
A. Bendale and T. Boult.
\newblock Towards open set deep networks.
\newblock In \emph{CVPR}, 2016.

\bibitem{rudd2018}
E. Rudd, L. Jain, W. Scheirer, and T. Boult.
\newblock The extreme value machine.
\newblock \emph{IEEE TPAMI}, 2018.

\bibitem{masud2011}
M. Masud et al.
\newblock Classification and novel class detection in concept-drifting data
streams under time constraints.
\newblock \emph{IEEE TKDE}, 2011.

\bibitem{haque2016sand}
A. Haque et al.
\newblock SAND: Semi-supervised adaptive novel class detection and
classification over data stream.
\newblock In \emph{AAAI}, 2016.

\bibitem{haque2016echo}
A. Haque et al.
\newblock ECHO: A data stream classification method for detecting novel class
drift.
\newblock In \emph{ICDE}, 2016.

\bibitem{faria2016}
E. de Faria, A. Gama, and A. Carvalho.
\newblock Novelty detection algorithm for data streams multi-class problems.
\newblock \emph{Data Mining and Knowledge Discovery}, 2016.

\bibitem{mu2017}
X. Mu, F. Zhu, J. Du, E.-P. Lim, and Z.-H. Zhou.
\newblock Classification under streaming emerging new classes.
\newblock \emph{IEEE TKDE}, 2017.

\bibitem{vaze2022}
S. Vaze et al.
\newblock Generalized category discovery.
\newblock In \emph{CVPR}, 2022.

\bibitem{cao2022}
K. Cao, M. Brbi\'c, and J. Leskovec.
\newblock Open-world semi-supervised learning.
\newblock In \emph{ICLR}, 2022.

\bibitem{zhang2022}
Z. Zhang et al.
\newblock Grow and Merge: A unified framework for continuous categories
discovery.
\newblock In \emph{NeurIPS}, 2022.

\bibitem{zhao2023}
B. Zhao and O. Mac Aodha.
\newblock Incremental generalized category discovery.
\newblock In \emph{ICCV}, 2023.

\bibitem{aljundi2017}
R. Aljundi, P. Chakravarty, and T. Tuytelaars.
\newblock Expert Gate: Lifelong learning with a network of experts.
\newblock In \emph{CVPR}, 2017.

\bibitem{gururangan2021}
S. Gururangan et al.
\newblock DEMix layers: Disentangling domains for modular language modeling.
\newblock arXiv:2108.05036, 2021.

\bibitem{li2022}
M. Li et al.
\newblock Branch-Train-Merge: Embarrassingly parallel training of expert
language models.
\newblock arXiv:2208.03306, 2022.

\bibitem{huang2024}
C. Huang et al.
\newblock LoraHub: Efficient cross-task generalization via dynamic LoRA
composition.
\newblock In \emph{COLM}, 2024.

\bibitem{hendrycks2017}
D. Hendrycks and K. Gimpel.
\newblock A baseline for detecting misclassified and out-of-distribution
examples in neural networks.
\newblock In \emph{ICLR}, 2017.

\bibitem{lee2018}
K. Lee et al.
\newblock A simple unified framework for detecting out-of-distribution samples
and adversarial attacks.
\newblock In \emph{NeurIPS}, 2018.

\bibitem{mukhoti2023}
J. Mukhoti et al.
\newblock Deep deterministic uncertainty: A new simple baseline.
\newblock \emph{TMLR}, 2023.

\bibitem{sun2022}
Y. Sun et al.
\newblock Out-of-distribution detection with deep nearest neighbors.
\newblock In \emph{ICML}, 2022.

\bibitem{vanamersfoort2020}
J. van Amersfoort et al.
\newblock Uncertainty estimation using a single deep deterministic neural
network.
\newblock In \emph{ICML}, 2020.

\bibitem{liu2020}
J. Liu et al.
\newblock Simple and principled uncertainty estimation with deterministic deep
learning via distance awareness.
\newblock In \emph{NeurIPS}, 2020.

\bibitem{ester1996}
M. Ester, H.-P. Kriegel, J. Sander, and X. Xu.
\newblock A density-based algorithm for discovering clusters in large spatial
databases with noise.
\newblock In \emph{KDD}, 1996.

\bibitem{campello2013}
R. Campello, D. Moulavi, and J. Sander.
\newblock Density-based clustering based on hierarchical density estimates.
\newblock In \emph{PAKDD}, 2013.

\bibitem{sarfraz2019}
S. Sarfraz, V. Sharma, and R. Stiefelhagen.
\newblock Efficient parameter-free clustering using first neighbor relations.
\newblock In \emph{CVPR}, 2019.

\bibitem{requicha1980}
A. Requicha.
\newblock Representations for rigid solids: Theory, methods, and systems.
\newblock \emph{ACM Computing Surveys}, 1980.

\bibitem{park2019}
J. Park et al.
\newblock DeepSDF: Learning continuous signed distance functions for shape
representation.
\newblock In \emph{CVPR}, 2019.

\bibitem{tipping1999}
M. Tipping and C. Bishop.
\newblock Probabilistic principal component analysis.
\newblock \emph{JRSS B}, 1999.

\bibitem{ghahramani1996}
Z. Ghahramani and G. Hinton.
\newblock The EM algorithm for mixtures of factor analyzers.
\newblock Technical Report CRG-TR-96-1, University of Toronto, 1996.

\bibitem{oquab2024}
M. Oquab et al.
\newblock DINOv2: Learning robust visual features without supervision.
\newblock \emph{TMLR}, 2024.

\end{thebibliography}

\appendix

\section{Final Branch Disposition}

\begin{table}[h]
\centering
\caption{Registered v8 milestone disposition.}
\begin{tabular}{lll}
\toprule
Milestone & Decision & Final status \\
\midrule
M45 & Endpoint, schemas, replay & Passed \\
M46 & Threshold and interface audit & Passed \\
M47 & Equal-budget utility selection & Outcome D \\
M48 & Learned selection & Blocked by M47 \\
M49 & Two local-residual attempts & No residual retained \\
M50/E12 & Untouched lifecycle qualification & Blocked by M47 \\
\bottomrule
\end{tabular}
\end{table}

\section{Artifact Verification}

The final replay reports:

\begin{itemize}
  \item four verified milestone indexes;
  \item 11 verified indexed artifacts;
  \item six reproduced conclusion operands;
  \item byte-identical output across two executions;
  \item no training-data loading; and
  \item no final-label access.
\end{itemize}

\section{Recent extension: the joint-budget scaling frontier}

Since the stage-wise study above, the programme measured the scaling
behaviour of the frozen patch-dictionary classifier. On a fixed DomainNet
subsample (138{,}000 train / 34{,}500 test rows, 345 classes, six domains),
the two-dimensional (dictionary-size $\times$ data) accuracy surface is
super-additive: scaling both axes jointly beats the sum of the single-axis
gains (excess +0.0097/+0.0101 over a +0.005 margin). The compute is
data-elastic --- head-width data-scaling steepness rises with dictionary
size --- and the atoms axis saturates near 6{,}144 atoms at full data. The
remaining gap to a frozen dense trunk is structural: a fixed low-rank code
geometry (effective rank $\approx$ 7.8) and a thin positive-margin argmax
tail. All comparisons are cost-matched; no dominance claim is made. A
self-contained report with reproduction steps is available at
\texttt{analysis/RESEARCH\_REPORT\_v19.md}; the registered literature
surveys (HTN-style routing; programmatic primitives and hybrid routers) are
documented in \texttt{analysis/HTN\_ROUTING\_LITERATURE\_REVIEW.md} and
\texttt{analysis/PROGRAMMATIC\_PRIMITIVES\_LITERATURE\_REVIEW.md}.

\end{document}
